% \usepackage{booktabs,threeparttable} required

\begin{table}[t]
\centering
\begin{threeparttable}
\caption{Stage 2 ILP solve time vs. requested pool size $\alpha{\times}k$ ($\alpha=10$ fixed,
\texttt{top\_n}=5000, $\sigma=0.60$), tier\_10k.}
\label{tab:wdc-ilp-scaling-tier10k}
\begin{tabular}{rrrrrr}
\toprule
$k$ & $\alpha{\times}k$ & Queries\tnote{a} & Falls short\tnote{b} & Mean ILP time (ms)\tnote{c} & Feasible\tnote{d} \\
\midrule
 10 &  100 & 30 & 0 &  2.33 & 22/30 \\
 20 &  200 & 30 & 0 &  1.49 & 16/30 \\
 30 &  300 & 30 & 0 &  1.58 & 15/30 \\
 40 &  400 & 30 & 1 &  2.22 & 14/30 \\
 50 &  500 & 30 & 1 &  1.93 & 11/30 \\
 60 &  600 & 30 & 2 &  1.83 & 11/30 \\
 70 &  700 & 30 & 2 &  6.30 & 11/30 \\
 80 &  800 & 30 & 3 &  3.75 & 11/30 \\
 90 &  900 & 30 & 4 &  4.51 & 11/30 \\
100 & 1000 & 30 & 5 & 37.65\tnote{e} & 10/30 \\
170\tnote{f} & 1700 & 30 & 19 & 6.64 & 9/30 \\
\bottomrule
\end{tabular}
\begin{tablenotes}
\footnotesize
\item[a] Queries that survived the pre-Stage-1 filter ($|D_{union}| \geq \max(k)$ within that row's
  sweep -- $\max(k)=100$ for the $k$=10..100 rows above, run together as one sweep; the $k$=170 row
  is its own singleton run, so its filter uses $\max(k)=170$ instead), out of 30 selected -- 0
  dropped at every $k$ shown here.
\item[b] "Falls short": queries where $\alpha{\times}k \geq |D_{union}|$ for that query -- Stage 1
  could not select a genuine $\alpha{\times}k$ pool and fell back to $P=D$ entirely (no Dinkelbach
  run for that query at that $k$). Grows with $k$ as the requested pool grows past what retrieval
  actually returned; Stage 2 still runs on the smaller, clamped $P$ and its time is still included
  in column [c].
\item[c] Mean Stage 2 (\texttt{scipy.optimize.milp}) solve time over all 30 queries at this $k$,
  including clamped cells.
\item[d] Number of the 30 queries where Stage 2 returned a feasible solution ($F_R \geq \tau$
  achievable within $P$), out of 30.
\item[e] \textbf{Single-query outlier, not a real trend.} Investigated by re-running this exact
  cell (tier\_10k, $k$=100, same 30 queries, same $D_{union}$ per query -- retrieval and query
  selection don't depend on $k$) three times. In every rerun, 29 of the 30 queries solved Stage 2
  in 1--7\,ms, consistent with every other $k$ in this table. One specific query --
  \texttt{1\_1438042987127.36\_20150728002307-00182-ip-10-236-191-2\_150862475\_5.csv}, the one
  case in this cell with a full, unclamped $n_P=1000$ pool and \texttt{feasible=True} (the largest
  and only fully-populated ILP instance among the 30) -- dominated the mean every time, but its OWN
  wall-clock solve time was itself unstable across reruns: 684\,ms, 684\,ms, then 126\,ms on a
  third rerun of the identical instance. Excluding it, the other 29 queries' mean is
  $\sim$2.9\,ms, in line with every other $k$ row. This points to branch-and-bound search-path
  sensitivity (system/thread-scheduling-dependent solve order in HiGHS, not a correctness issue --
  the C4 exactness/post-solve-verification invariants that \texttt{stage2\_ilp.py} enforces are
  about the RESULT, never wall-clock time) on this one genuinely hard, near-full-pool instance, not
  a real "$k$=100 is slow" scaling effect. The reported 37.65\,ms mean should be read as
  ``dominated by one unstable outlier out of 30,'' not as a representative per-query cost at this
  pool size.
\item[f] \textbf{Max-$k$ row.} $k$ chosen so $\alpha{\times}k$ reaches the mean $|D_{union}|$
  actually retrieved at this tier/\texttt{top\_n}/$\sigma$ (1702, measured fresh for this row --
  supersedes the 1643 figure in \texttt{topn\_vs\_dunion\_tables.tex}, which predates the
  categorical-only HNSW fix and is now stale). $|D_{union}|$ is bimodal across the 30 queries here
  (26 queries in $[333,1899]$, 4 in $[3376,4120]$), so a $k$ sized to the MEAN clamps most of the
  smaller-yield queries (19 of 30, column [b]) while leaving the few large-yield queries a genuine,
  non-clamped Dinkelbach selection -- there is no single $k$ at which $\alpha{\times}k$ reaches the
  mean AND leaves every query unclamped, given this split distribution.
\end{tablenotes}
\end{threeparttable}
\end{table}

\begin{table}[t]
\centering
\begin{threeparttable}
\caption{Stage 2 ILP solve time vs. requested pool size $\alpha{\times}k$ ($\alpha=10$ fixed,
\texttt{top\_n}=5000, $\sigma=0.60$), tier\_100k.}
\label{tab:wdc-ilp-scaling-tier100k}
\begin{tabular}{rrrrrr}
\toprule
$k$ & $\alpha{\times}k$ & Queries\tnote{a} & Falls short\tnote{b} & Mean ILP time (ms)\tnote{c} & Feasible\tnote{d} \\
\midrule
 10 &  100 & 30 & 0 & 2.50 & 20/30 \\
 20 &  200 & 30 & 0 & 2.06 & 13/30 \\
 30 &  300 & 30 & 0 & 1.91 & 13/30 \\
 40 &  400 & 30 & 0 & 2.49 & 12/30 \\
 50 &  500 & 30 & 1 & 4.34 & 11/30 \\
 60 &  600 & 30 & 1 & 3.90 & 10/30 \\
 70 &  700 & 30 & 2 & 8.20 & 10/30 \\
 80 &  800 & 30 & 2 & 2.44 & 10/30 \\
 90 &  900 & 30 & 2 & 5.15 & 10/30 \\
100 & 1000 & 30 & 3 & 3.34 &  9/30 \\
639\tnote{f} & 6390 & 29\tnote{g} & 23 & 41.88 & 7/29 \\
\bottomrule
\end{tabular}
\begin{tablenotes}
\footnotesize
\item[a] Queries that survived the pre-Stage-1 filter ($|D_{union}| \geq \max(k)$ within that row's
  sweep -- $\max(k)=100$ for the $k$=10..100 rows above, run together as one sweep; the $k$=639 row
  is its own singleton run, so its filter uses $\max(k)=639$ instead), out of 30 selected.
\item[b] "Falls short": queries where $\alpha{\times}k \geq |D_{union}|$ for that query -- Stage 1
  could not select a genuine $\alpha{\times}k$ pool and fell back to $P=D$ entirely (no Dinkelbach
  run for that query at that $k$). Grows with $k$ as the requested pool grows past what retrieval
  actually returned; Stage 2 still runs on the smaller, clamped $P$ and its time is still included
  in column [c]. Fewer queries fall short here than at tier\_10k despite the identical $\alpha,k$
  sweep, because tier\_100k's $|D_{union}|$ at \texttt{top\_n}=5000 is roughly 4x larger on average
  (mean 6387 vs. 1702 at tier\_10k, both measured fresh for this table).
\item[c] Mean Stage 2 (\texttt{scipy.optimize.milp}) solve time over all 30 (or 29, $k$=639 row)
  queries at this $k$, including clamped cells.
\item[d] Number of queries where Stage 2 returned a feasible solution ($F_R \geq \tau$ achievable
  within $P$), out of the queries used at that $k$ (column [a]).
\item[f] \textbf{Max-$k$ row.} $k$ chosen so $\alpha{\times}k$ reaches the mean $|D_{union}|$
  actually retrieved at this tier/\texttt{top\_n}/$\sigma$ (6387, measured fresh for this row --
  supersedes the 6327 figure in \texttt{topn\_vs\_dunion\_tables.tex}, which predates the
  categorical-only HNSW fix and is now stale). $|D_{union}|$ is bimodal across the 30 queries here
  (24 queries in $[406,2044]$, 6 in $[23277,30750]$), so a $k$ sized to the MEAN clamps essentially
  every smaller-yield query (23 of the 29 that reach Stage 1, column [b]) while leaving the 6
  large-yield queries a genuine, non-clamped Dinkelbach selection -- there is no single $k$ at which
  $\alpha{\times}k$ reaches the mean AND leaves every query unclamped, given this split
  distribution.
\item[g] One of the 30 selected queries ($|D_{union}|=406$) falls below this row's own pre-Stage-1
  threshold ($k$=639, see note [a]) and is dropped before Stage 1 runs -- 29 queries remain.
\end{tablenotes}
\end{threeparttable}
\end{table}
